\chapter{Smart home energy management}\label{SmartEnergyStatement}
\chapterprecis{%
  This chapter is from DD2520 Applied Crypto, spring 2026, at KTH\@.
}

From 1 January 2027, all electricity providers in Sweden are required to
charge an \emph{effect tariff} (due to EU regulations).
This means that the electricity bill will be based on the price of electricity
at the time of use, but also on the maximum power (effect) used during the
billing period.
This is done to encourage users to distribute their electricity use over time,
so that they don't run the oven, charge their electric car, run the dishwasher
and laundry machine all at the same time.

To help households distribute their electricity use, we want smart products to
coordinate when they consume power.
For instance, the car charger and the heating system can automatically lower
their consumption when they detect that the oven, vacuum cleaner, dishwasher
or laundry machine are running.
The dishwasher can delay its start until power consumption is lower or the
price is cheaper.

\ltnote{%
  The idea is to use more open form of pQBL~\parencite{pQBL}.
  This means that we start with a question.
  This lets the reader think about the problem, try to explore it based on
  what they currently know.
  When we later tell them, the knowledge will be more
  reinforced~\parencite{Szekely1950}.
  This also helps the reader to discern relevant
  aspects~\parencite{NecessaryConditionsOfLearning}.

  The intention of \cref{EnergyBenefitsDrawbacks,EnergyAlleviateDrawbacks} is
  to let the reader discover the relevant security and privacy properties that
  we want to capture in the coming protocol design.
}
\begin{exercise}\label{EnergyBenefitsDrawbacks}
  What are the benefits and drawbacks of having smart devices coordinate
  their electricity use?
\end{exercise}

Benefits include lower electricity bills by avoiding peak-effect charges,
reduced strain on the electricity grid, and automatic comfort optimization
without manual intervention.

Drawbacks include the risk of exposing private habits through power-usage
data (when you cook, when you're home, when you sleep), the risk of
unauthorized control of devices (an attacker turning off your heating in
winter), and dependence on the security of wireless communication.

\begin{exercise}\label{EnergyAlleviateDrawbacks}
  What can we do to alleviate the drawbacks while enjoying the benefits?
\end{exercise}

We can use cryptography to protect the communication between devices:
encryption for confidentiality, authentication to prevent unauthorized
commands, and integrity protection to detect tampering.
We must also carefully consider what data the electricity provider needs and
minimize what we share.

\begin{exercise}
  Think about how you would design a system using cryptographic primitives to
  ensure the security and privacy of the smart energy system.
\end{exercise}

\ltnote{%
  The idea is to create the simplest solution possible.
  Then we will focus one aspect at a time by making one small change to the
  solution at a time.
  This will introduce the contrast necessary to bring the aspect into
  focus \parencite[as dictated by][]{NecessaryConditionsOfLearning}.
  We will do this until we have touched on all aspects and arrived at the
  perfect solution.
  (One could also think of this as a sequence of games, as used in
  cryptography \parencite{Shoup2004Sequences}.
  Arguably, the reason why the sequences of games method by
  \citeauthor{Shoup2004Sequences} is good for humans is explained by
  \citeauthor{NecessaryConditionsOfLearning}.)
}

In this chapter, we will iteratively develop a cryptographic solution,
refining it step by step.
\chapterminitoc


\section{A first solution}\label{EnergyFirstSolution}

We have the following entities from \cref{SmartEnergyStatement}:
\begin{itemize}
  \item the power meter~\(M\), which measures the household's current power
    consumption and reports the maximum to the electricity provider,
  \item the controller~\(H\) (home hub), which coordinates electricity use
    among the devices,
  \item the smart devices~\(D_1, D_2, \ldots, D_n\), such as dishwashers,
    car chargers, heating systems and ovens,
  \item the electricity provider~\(E\), who bills the household based on
    consumption and peak effect.
\end{itemize}

\begin{exercise}
  What is the goal of the system, in terms of these entities?
\end{exercise}

The goal is that the controller~\(H\) receives power-usage data from the
meter~\(M\) and from each device~\(D_i\), and sends commands to the
devices to adjust their consumption.
The provider~\(E\) receives billing data from the meter~\(M\).
All communication is wireless (radio), which means anyone nearby can
eavesdrop or inject messages.

\begin{exercise}
  Design a cryptographic system that ensures the confidentiality, integrity
  and authenticity of the communication between these entities.
\end{exercise}

The simplest design uses a shared symmetric key between the controller~\(H\)
and each device~\(D_i\).
The controller also shares a key with the meter~\(M\).

\subsection{Keys and communication}

We have a shared symmetric key~\(\KHD\) between the controller~\(H\) and
each device~\(D_i\), and a key~\(\KHM\) between the controller and the
meter~\(M\).

\begin{exercise}
  How should these keys be stored and managed?
\end{exercise}

Each device stores its key in a secure element (\eg a small TPM or secure
microcontroller).
The controller stores all keys and must be physically protected---it is
typically installed in a locked utility cupboard or similar.

The keys can be pre-provisioned during manufacturing: each device comes with
a unique key, and the controller learns the key during a pairing procedure
(\eg by scanning a QR code on the device or by entering a code printed on
the device's label).

\subsection{A simple command protocol}

We derive an encryption key~\(\KEHD\) and a MAC key~\(\KMHD\) from the
shared key~\(\KHD\):
\begin{align*}
  \KEHD &= \KDF[\Enc, \KHD], \\
  \KMHD &= \KDF[\MAC, \KHD].
\end{align*}

When the controller~\(H\) wants to send a command to device~\(D_i\)
(\eg \enquote{suspend until 14:00}), it encrypts and MACs the message:
\begin{align*}
  c &= \Enc[\KEHD, m], \\
  t &= \MAC[\KMHD, c].
\end{align*}

The device verifies the MAC before decrypting.
Similarly, when a device reports its power consumption to the controller, it
encrypts and MACs using the same keys.

See \cref{fig:EnergyFirstProtocol} for the protocol.

\begin{figure}
  \begin{sidecaption}[Command protocol between \(H\) and \(D_i\)]{%
    The controller~\(H\) sends a command to device~\(D_i\).
    The message is encrypted for confidentiality and MACed for integrity
    and authentication.
    The device verifies the MAC before acting on the command.
  }[fig:EnergyFirstProtocol]
  \flushright
  \begin{msc}
    [/msc/action width=0.4\linewidth]
    {Command from \(H\) to \(D_i\)}
      \declinst{ctrl}{Controller \(H\)}{\(\KHD\)}
      \declinst{dev}{Device \(D_i\)}{\(\KHD\)}
      \mscaction{%
        \(c = \Enc[\KEHD, m]\),
        \(t = \MAC[\KMHD, c]\)
      }{ctrl}
      \nextlevel[4]
      \mess{\(c, t\)}{ctrl}{dev}
      \nextlevel
      \mscaction{Verify \(t\), decrypt \(c\)}{dev}
      \mscaction{Execute command}{dev}
      \nextlevel
  \end{msc}
  \end{sidecaption}
\end{figure}

\subsection{Consequences of the design}\label{EnergyConsequencesFirst}

\begin{exercise}
  Evaluate the design.
  What are its consequences (advantages and disadvantages) that we must
  consider?
\end{exercise}

The design has several consequences:
\begin{enumerate}
  \item Simple and efficient: symmetric cryptography is fast, which suits
    resource-constrained devices like light bulbs and appliances.
  \item Key management is manual: the user must pair each device individually
    with the controller.
    Adding or replacing a device requires physical interaction.
  \item No replay protection: an attacker who records a valid
    \enquote{suspend} command can replay it later to disrupt the household.
  \item Privacy is not addressed: the meter~\(M\) and the provider~\(E\)
    learn detailed consumption patterns.
  \item If any device is compromised, the attacker learns the key~\(\KHD\)
    and can impersonate the controller to that device.
    But other devices are not affected (different keys).
\end{enumerate}


\section{Protecting against replay attacks}\label{EnergyReplayProtection}

\ltnote{%
  We introduce replay protection to create contrast with the previous design
  where messages could be replayed.
  Students should discern that encryption and authentication alone do not
  prevent replay---a separate mechanism (sequence numbers, timestamps, or
  nonces) is needed.
  The variation pattern: we keep the encryption and MAC invariant while
  adding replay protection.
}

In the first design, an attacker can record a valid encrypted and MACed
command and replay it later.
For example, replaying a \enquote{suspend charging} command to the car
charger every evening could prevent the car from ever charging.

\begin{exercise}
  How can we prevent replay attacks in this protocol?
\end{exercise}

\subsection{Sequence numbers}

We add a monotonically increasing sequence number~\(s\) to each message.
The sender increments~\(s\) for each message sent.
The receiver keeps track of the highest sequence number received and rejects
any message with a sequence number that is not strictly greater.

The MAC now covers both the ciphertext and the sequence number:
\begin{align*}
  c &= \Enc[\KEHD, m \concat s], \\
  t &= \MAC[\KMHD, c \concat s].
\end{align*}

Alternatively, we can use timestamps if the devices have synchronized clocks.
However, clock synchronization in resource-constrained devices is
non-trivial and introduces another attack surface (what if an attacker
manipulates the time source?).
Sequence numbers are simpler and don't require synchronized clocks.

\begin{exercise}
  What happens if the controller or a device loses power and resets its
  sequence counter?
\end{exercise}

If a device resets, its sequence counter returns to zero.
The controller would then reject its messages (sequence number too low).
To handle this, the sequence counter should be stored in non-volatile memory.
Alternatively, the device and controller can run a re-synchronization
protocol after a reset, but this must be authenticated to prevent an
attacker from forcing a counter reset.

\subsection{Consequences of the design}

\begin{enumerate}
  \item Replay attacks are prevented: old messages are rejected.
  \item Sequence counters must be stored in non-volatile memory to survive
    power loss.
  \item Out-of-order delivery is not supported: if messages arrive out of
    order, they will be rejected.
    This is acceptable for command-and-control protocols but may need a
    sliding window for high-throughput telemetry.
  \item The remaining issues from \cref{EnergyConsequencesFirst} still
    apply.
\end{enumerate}


\section{Device authentication and pairing}\label{EnergyPairing}

\ltnote{%
  We now address the question of how devices join the network securely.
  This creates contrast with the previous assumption that keys are
  pre-provisioned.
  Students should discern that secure pairing is a bootstrapping problem:
  how do you establish a shared secret when you have no prior shared secret?
  The variation pattern: we keep the communication protocol invariant while
  varying the key establishment method.
}

In \cref{EnergyFirstSolution}, we assumed that each device is pre-provisioned
with a key known to the controller.
In practice, a user buys a new smart light bulb and wants to add it to their
home network.
The bulb has no prior relationship with the controller.

\begin{exercise}
  How can a new device securely join the home network and establish a shared
  key with the controller?
\end{exercise}

\subsection{Out-of-band pairing}

The standard approach is \emph{out-of-band} (OOB) pairing: using a separate,
trusted channel to bootstrap the key exchange.
Common methods:
\begin{description}
  \item[QR code] The device has a QR code (printed on its body or packaging)
    containing a pre-shared secret or a public key.
    The user scans it with the controller's app.
  \item[NFC tap] The user holds the device near the controller (or a phone
    acting as the controller's interface) to exchange pairing data via
    near-field communication.
  \item[PIN entry] The device displays or is labelled with a PIN code.
    The user enters it into the controller.
\end{description}

After the OOB exchange provides a shared secret (or authenticated public
keys), the devices can run a key-agreement protocol (\eg Diffie--Hellman)
over the wireless channel, authenticated by the OOB secret, to establish
the session key~\(\KHD\).

\begin{exercise}
  What are the security properties of each pairing method?
  Which is most suitable for the different types of devices in our system?
\end{exercise}

\begin{description}
  \item[QR code] Works well for devices that are too simple for a display or
    keypad (\eg light bulbs).
    The code is static, so if an attacker photographs it, they can pair
    later.
    The code should be inside the packaging and treated as confidential
    until pairing.
  \item[NFC tap] Requires proximity (a few centimetres), which provides
    physical authentication.
    Good for devices the user can physically approach.
    Harder to intercept than a QR code.
  \item[PIN entry] Requires a display or label on the device and an input
    method on the controller.
    The PIN must be long enough to resist brute-force attempts (at least
    6 digits).
    Suitable for devices with displays (\eg ovens, car chargers).
\end{description}

\subsection{Certificate-based onboarding}

For a more scalable approach, especially in larger installations, we can use
certificates.
Each device is manufactured with a device certificate signed by the
manufacturer's CA (Certificate Authority).
The controller~\(H\) has its own key pair~\((\SKH, \VKH)\), certified by a
home CA or self-signed.

During pairing:
\begin{enumerate}
  \item The device presents its manufacturer certificate.
  \item The controller verifies the certificate chain.
  \item They perform authenticated key exchange (using their respective
    keys) to establish~\(\KHD\).
\end{enumerate}

This approach doesn't require QR codes or PINs, but it requires the
controller to trust the manufacturer's CA\@.
It also assumes the manufacturer has not been compromised and that device
certificates have not been cloned.

\subsection{Consequences of the design}

\begin{exercise}
  Evaluate the pairing methods.
  What are the trade-offs?
\end{exercise}

\begin{description}
  \item[Usability vs.\ security] QR codes and NFC are easy for users but
    rely on physical security; certificate-based pairing is seamless but
    requires trust in manufacturers.
  \item[Scalability] Certificate-based pairing scales better to many
    devices, since no per-device manual step is needed.
  \item[Device cost] Certificates require more capable hardware (crypto
    coprocessor) in each device, increasing cost.
    Symmetric-key OOB pairing can work with very cheap hardware.
  \item[Revocation] If a device is compromised, we need a way to revoke its
    key.
    With pre-shared keys, the controller simply deletes the key.
    With certificates, a revocation mechanism is needed.
\end{description}


\section{Protecting wireless communication}\label{EnergyWirelessSecurity}

\ltnote{%
  We now focus on the wireless channel itself as an attack surface.
  This creates contrast with the previous focus on message-level security
  (encryption and MACs).
  Students should discern that even with authenticated and encrypted
  messages, the wireless medium introduces additional threats: jamming,
  traffic analysis, and physical-layer attacks.
  The variation pattern: we keep the cryptographic protocol invariant while
  varying the channel properties.
}

All communication in the system happens over wireless radio.
Our encrypted and authenticated messages protect against eavesdropping on
content and against message forgery.
But there are additional threats specific to wireless channels.

\begin{exercise}
  What threats exist at the wireless channel level, beyond what our
  encryption and authentication already address?
\end{exercise}

\subsection{Jamming}

An attacker can transmit noise on the radio frequency to prevent all
communication.
This is a denial-of-service attack.
If the devices cannot communicate, the controller cannot coordinate
power use, potentially causing the household to exceed its effect limit
and incur higher bills.

\begin{exercise}
  How can we mitigate jamming attacks?
\end{exercise}

Complete protection against jamming is difficult, but mitigations include:
\begin{description}
  \item[Frequency hopping] The devices agree (during pairing) on a
    pseudo-random frequency-hopping sequence.
    An attacker must jam all frequencies simultaneously, which requires
    much more power.
    This is standard in protocols like Bluetooth.
  \item[Graceful degradation] Devices should have sensible defaults when
    they cannot reach the controller.
    For example, the car charger could continue charging at a low rate
    rather than stopping entirely.
  \item[Detection and alerting] The controller can detect communication
    failures and alert the user.
\end{description}

\subsection{Traffic analysis}

Even with encryption, an attacker observing the wireless traffic can learn:
\begin{itemize}
  \item Which devices are communicating (by radio fingerprinting or
    observing communication patterns).
  \item When commands are sent (indicating changes in household activity).
  \item The volume of data, which may reveal whether devices are active.
\end{itemize}

\begin{exercise}
  How can we mitigate traffic analysis?
\end{exercise}

Mitigations include:
\begin{itemize}
  \item \emph{Constant-rate traffic}: Devices transmit at a fixed rate,
    sending dummy messages when there is no real data.
    This hides when actual commands are sent.
  \item \emph{Uniform message sizes}: Pad all messages to the same length
    so that message type cannot be inferred from size.
  \item \emph{Randomized timing}: Add random delays to message
    transmission to obscure patterns.
\end{itemize}

These measures increase power consumption and bandwidth use, so they must
be weighed against the limited resources of battery-powered devices.

\subsection{Consequences of the design}

\begin{description}
  \item[Availability] Frequency hopping and graceful degradation improve
    resilience against jamming.
  \item[Privacy] Traffic-analysis countermeasures protect activity patterns
    but consume more energy and bandwidth.
  \item[Complexity] Frequency hopping requires shared state between devices
    and controller; constant-rate traffic requires careful timing management.
  \item[Battery life] Continuous dummy traffic significantly reduces battery
    life for battery-powered devices (sensors, smart plugs).
    A reasonable trade-off may be to use constant-rate traffic only for
    sensitive devices.
\end{description}


\section{Privacy of energy data}\label{EnergyPrivacy}

\ltnote{%
  We now address the privacy dimension of the system.
  This creates contrast with the previous focus on securing communication
  between the controller and devices.
  Students should discern that the electricity provider~\(E\) learning
  detailed consumption data is a separate threat from a wireless
  eavesdropper.
  The variation pattern: we keep the local protocol invariant while varying
  the data shared with external parties.
}

The power meter~\(M\) records the household's power consumption over time.
The electricity provider~\(E\) needs this data for billing.
But detailed consumption data reveals intimate information about the
household's daily life.

\begin{exercise}
  What can an electricity provider (or an attacker who compromises the
  meter data) learn from detailed power consumption data?
\end{exercise}

From fine-grained power data, one can infer:
\begin{itemize}
  \item When residents are home and when they are away.
  \item When they wake up and go to sleep.
  \item What appliances they use (each appliance has a characteristic power
    signature).
  \item Whether they have an electric car and when they charge it.
  \item How many people live in the household (from shower and cooking
    patterns).
\end{itemize}

This is called \emph{non-intrusive load monitoring} (NILM) and is a
well-studied technique.

\begin{exercise}
  How can we design the system to minimize the privacy impact while still
  enabling the provider to bill correctly?
\end{exercise}

\subsection{Data minimization}

The principle of \emph{data minimization} (central to GDPR) states that we
should collect only the data strictly necessary for the purpose.

For billing under effect tariffs, the provider needs:
\begin{itemize}
  \item Total energy consumed during the billing period.
  \item Maximum power (effect) observed during the billing period.
\end{itemize}

The provider does \emph{not} need minute-by-minute or second-by-second
consumption readings.
Therefore, the meter should report only aggregated values.

\subsection{Aggregation at the meter}

The meter~\(M\) can compute the aggregated billing data locally and report
only the result to the provider~\(E\):
\begin{itemize}
  \item Total consumption: \(\sum_{t} P(t) \cdot \Delta t\).
  \item Peak effect: \(\max_{t} P(t)\) over the billing period.
\end{itemize}

The meter sends these values to the provider, encrypted and signed, at the
end of each billing period.
The detailed time series stays within the household.

\begin{exercise}
  The controller~\(H\) needs detailed data from the meter for coordination.
  How do we handle this without leaking data to the provider?
\end{exercise}

The meter~\(M\) shares detailed readings with the controller~\(H\) over the
local wireless network, protected by the key~\(\KHM\).
The provider only receives the aggregated billing data.
This separation is enforced by the meter's software, which must be trusted.

\subsection{Consequences of the design}

\begin{exercise}
  Evaluate the privacy design.
  What are its strengths and weaknesses?
\end{exercise}

\begin{description}
  \item[Privacy] The provider learns only aggregate billing data, not
    detailed usage patterns.
    This is a significant improvement over reporting fine-grained data.
  \item[Trust in the meter] The meter must be trusted to correctly compute
    aggregates and not leak detailed data.
    This is a single point of trust.
  \item[Provider's perspective] The provider cannot verify that the meter
    is reporting honestly.
    We address this in \cref{EnergyMeteringIntegrity}.
  \item[Local privacy] Within the household, the controller sees detailed
    data.
    If the controller is cloud-connected (\eg for remote access), this
    data could leak to the cloud provider.
\end{description}


\section{Integrity of metering data}\label{EnergyMeteringIntegrity}

\ltnote{%
  We now address the tension between privacy (minimizing data shared with
  the provider) and integrity (the provider needs trustworthy billing data).
  This creates contrast with the previous section where we trusted the
  meter.
  Students should discern that data minimization and verifiability are in
  tension and that hardware-based trust can bridge this gap.
  The variation pattern: we keep the privacy requirement invariant while
  adding integrity assurance for the provider.
}

The electricity provider~\(E\) needs to trust that the billing data from the
meter~\(M\) is correct.
A dishonest household could tamper with the meter to report lower consumption
or a lower peak effect, reducing their bill.

\begin{exercise}
  How can the provider verify that the meter's reports are correct?
\end{exercise}

\subsection{Signed meter readings}

Each meter~\(M\) has a signing key pair~\((\SKM, \VKM)\), provisioned during
manufacturing.
The public key~\(\VKM\) is registered with the provider~\(E\) (possibly via
the manufacturer's certificate chain).

The meter signs each billing report:
\[
  \sigma = \Sign[\SKM, (\text{meter\_id}, \text{period}, \text{total\_kWh},
  \text{peak\_kW}, \text{timestamp})].
\]

The provider verifies:
\(\Ver[\VKM, \text{report}, \sigma] \stackrel{?}{=} \text{true}\).

This ensures that the report was produced by the authentic meter and has not
been modified in transit.

\begin{exercise}
  The signature proves the report came from the meter, but what if the
  meter's hardware has been physically tampered with?
\end{exercise}

\subsection{Tamper-evident meters}

Physical security measures complement the cryptographic ones:
\begin{description}
  \item[Tamper-evident seals] Visible indicators that the meter's enclosure
    has been opened.
  \item[Tamper-responsive hardware] The meter's secure element erases the
    signing key if physical intrusion is detected, making further signed
    reports impossible.
  \item[Secure boot] The meter verifies its own firmware integrity at
    start-up, refusing to operate if the software has been modified.
  \item[Remote attestation] The meter can prove to the provider that it is
    running genuine, unmodified firmware.
\end{description}

\subsection{The communication channel to the provider}

\begin{exercise}
  How should the meter communicate with the electricity provider?
\end{exercise}

The meter can communicate with the provider through several channels:
\begin{itemize}
  \item \emph{Power-line communication (PLC)}: Data transmitted over the
    existing electrical wiring.
    Common in smart metering deployments.
  \item \emph{Cellular (4G/5G)}: The meter has its own SIM card and
    communicates directly with the provider.
  \item \emph{Via the household's internet connection}: Cheaper but relies
    on the household's network, which the household controls.
\end{itemize}

Regardless of the channel, the meter authenticates to the provider and
establishes an encrypted connection (\eg TLS with mutual authentication).
The signed billing reports provide an additional layer of integrity that
holds even if the transport channel is compromised.

\subsection{Consequences of the design}

\begin{description}
  \item[Billing integrity] Signed reports provide non-repudiation: the
    household cannot claim the meter reported incorrectly (assuming the
    meter hardware is intact).
  \item[Privacy preserved] The provider still receives only aggregate data;
    the signature doesn't require sharing more information.
  \item[Hardware cost] Tamper-evident meters with secure elements are more
    expensive than simple meters.
  \item[Key management] The provider must maintain a database of all meter
    public keys and manage certificate revocation for decommissioned meters.
  \item[Single device trust] The entire billing integrity depends on the
    meter's hardware and software being trustworthy.
    This is a strong assumption, but it is standard practice in metering.
\end{description}


\section{Firmware updates and device lifecycle}\label{EnergyFirmware}

\ltnote{%
  We introduce firmware security to create contrast with the
  protocol-level security we have focused on so far.
  Students should discern that even correct protocols can be undermined if
  the device software is modified by an attacker.
  The variation pattern: we keep the protocol invariant while varying the
  trust in device software.
}

Smart devices, the controller, and the meter all run software (firmware) that
implements our protocols.
This firmware will need updates to fix bugs, patch security vulnerabilities,
and add features.

\begin{exercise}
  What threats are associated with firmware updates?
  How can an attacker exploit the update mechanism?
\end{exercise}

An attacker who can install malicious firmware on a device can:
\begin{itemize}
  \item Extract cryptographic keys from the device.
  \item Disable security checks (MAC verification, encryption).
  \item Cause the device to malfunction (\eg the heating system running
    at maximum in summer).
  \item Use the device as a foothold to attack other devices on the
    network.
\end{itemize}

\subsection{Signed firmware updates}

Each manufacturer has a firmware-signing key pair
\((\SKManuf, \VKManuf)\).
The device stores the manufacturer's public key~\(\VKManuf\) in its secure
element.

Before installing a firmware update, the device verifies:
\[
  \Ver[\VKManuf, \text{firmware\_image}, \sigma_{\text{fw}}]
  \stackrel{?}{=} \text{true}.
\]

If verification fails, the update is rejected.

\subsection{Secure boot}

\begin{exercise}
  Even if we only install signed firmware, what happens if the device's
  storage is modified directly (by physical access)?
\end{exercise}

\emph{Secure boot} provides the answer: at every power-on, the device
verifies the integrity of its own firmware before executing it.
The chain of trust starts from an immutable boot ROM that verifies the
bootloader, which in turn verifies the main firmware.

If any stage fails verification, the device refuses to boot normally and
enters a recovery mode that only accepts signed firmware from the
manufacturer.

\subsection{Revoking compromised devices}

\begin{exercise}
  If a device is compromised (its keys extracted by an attacker), how
  should the system respond?
\end{exercise}

The controller should maintain a list of authorized devices.
When a device is reported compromised:
\begin{enumerate}
  \item The controller removes the device's key~\(\KHD\) from its
    database, refusing further communication.
  \item If certificate-based pairing is used, the device's certificate is
    added to a revocation list.
  \item The user is notified to replace or reset the device.
  \item If the compromised device's key was used for any shared state,
    affected keys should be rotated.
\end{enumerate}

\subsection{Consequences of the design}

\begin{description}
  \item[Software integrity] Signed firmware and secure boot prevent
    unauthorized software from running on devices.
  \item[Manufacturer trust] Devices must trust their manufacturer's
    signing key.
    If a manufacturer's key is compromised, all their devices are
    vulnerable.
  \item[Update delivery] Firmware updates must reach all devices reliably,
    including battery-powered devices that may be offline for extended
    periods.
  \item[Backward compatibility] Updated firmware must maintain protocol
    compatibility with devices that haven't been updated yet.
  \item[Device lifecycle] Eventually, manufacturers stop supporting old
    devices.
    There should be a mechanism to continue operating (with a clear
    security status) or to decommission them gracefully.
\end{description}


\section{Switching to authenticated encryption}\label{EnergyAEAD}

\ltnote{%
  We now replace the separate encrypt-then-MAC approach with AEAD to
  simplify the protocol and eliminate potential composition errors.
  This creates contrast with the earlier design using separate primitives.
  Students should discern that AEAD provides the same properties
  (confidentiality, integrity, authenticity) in a single, well-analyzed
  primitive, reducing the risk of subtle implementation mistakes.
  The variation pattern: we keep the security properties invariant while
  varying the cryptographic mechanism.
}

In our earlier design, we used separate encryption (CTR mode) and MAC
(HMAC) with two derived keys.
This \enquote{encrypt-then-MAC} composition is secure when done correctly,
but there are well-known pitfalls in combining separate primitives
(\eg MAC-then-encrypt, or forgetting to MAC the IV).

\begin{exercise}
  What is AEAD (Authenticated Encryption with Associated Data), and why
  might it be preferable to separate encrypt-then-MAC?
\end{exercise}

AEAD combines encryption and authentication into a single operation.
The most widely used AEAD construction is AES-GCM (Galois/Counter Mode),
which we also encountered in \cref{FirstSolution}.
Another modern option is ChaCha20-Poly1305, which is popular on
resource-constrained devices due to its software efficiency.

With AEAD, our protocol simplifies to:
\begin{align*}
  c &= \text{AEAD-Enc}[\KHD, \text{nonce}, m, \text{aad}],
\end{align*}
where the associated data (aad) can include the sequence number and device
identifier without encrypting them (they are authenticated but sent in
the clear).

\begin{exercise}
  What should go into the associated data?
\end{exercise}

The associated data should include:
\begin{itemize}
  \item The sequence number (for replay protection).
  \item The sender and receiver identifiers (to prevent messages intended
    for one device from being accepted by another).
  \item A protocol version number (to prevent cross-version attacks).
\end{itemize}

\subsection{Nonce management}

\begin{exercise}
  AEAD requires a unique nonce for each message under the same key.
  How should we manage nonces?
\end{exercise}

For AES-GCM, nonce reuse is catastrophic: it leaks the XOR of two
plaintexts and allows forgery of authentication tags.

We can use the sequence number as the nonce, since it is already guaranteed
to be unique and monotonically increasing.
This eliminates the need for a separate random nonce and avoids the birthday
bound problem for randomly generated nonces.

If using ChaCha20-Poly1305, the nonce is 96 bits and the same
sequence-number approach works well.

\subsection{Consequences of the design}

\begin{description}
  \item[Simplicity] A single primitive (AEAD) replaces two (encryption +
    MAC), reducing implementation complexity.
  \item[Fewer keys] We derive a single key per device pair instead of two
    (encryption + MAC).
  \item[Proven security] AEAD constructions are well-analyzed; the
    composition is guaranteed to be secure (unlike ad-hoc combinations).
  \item[Nonce sensitivity] AES-GCM requires strict nonce uniqueness.
    Sequence numbers provide this naturally.
  \item[Algorithm choice] AES-GCM benefits from hardware acceleration
    (AES-NI) on devices that have it.
    ChaCha20-Poly1305 is better for devices without hardware AES support.
\end{description}


\section{Legal and usability considerations}\label{EnergyLegalUsability}

\ltnote{%
  We now shift from cryptographic protocol design to legal and usability
  aspects.
  This creates contrast with the purely technical focus of previous
  sections.
  Students should discern that a technically secure system can still fail
  if it doesn't meet legal requirements or is too difficult for ordinary
  households to use.
  The variation pattern: we keep the technical design invariant while
  varying the evaluation perspective.
}

\begin{exercise}
  What legal requirements apply to a smart energy management system that
  processes household energy data?
\end{exercise}

\subsection{GDPR and energy data}

Under GDPR, energy consumption data is personal data because it can reveal
information about the household's inhabitants.
Key obligations:
\begin{description}
  \item[Purpose limitation] Data may only be processed for specified
    purposes (billing, grid management).
  \item[Data minimization] Collect only what is strictly necessary, as we
    designed in \cref{EnergyPrivacy}.
  \item[Storage limitation] Detailed data should not be retained longer
    than needed.
  \item[Data subject rights] Households can request access to their data,
    correction, or deletion.
  \item[Security] Appropriate technical measures (our cryptographic
    protocols) must protect the data.
\end{description}

\subsection{EU smart meter regulations}

The EU has specific regulations for smart metering:
\begin{itemize}
  \item Meters must be able to be read remotely by the provider.
  \item Consumers must have access to their own consumption data.
  \item Data must be provided in a standardized format.
  \item Consumers can choose to share their data with third parties (\eg
    energy-management services).
\end{itemize}

Our design must ensure that the consumer can share data selectively,
without giving the third party access to more data than necessary.

\subsection{Usability for non-technical households}

\begin{exercise}
  How do we make the system usable for households where nobody has
  technical expertise?
\end{exercise}

Key usability principles:
\begin{description}
  \item[Automatic coordination] Devices should coordinate without
    requiring the user to understand the underlying protocols.
    The system should work out of the box with sensible defaults.
  \item[Simple pairing] Adding a new device should require at most scanning
    a QR code or pressing a button.
    The user should not have to configure keys, algorithms, or
    security parameters.
  \item[Transparent operation] The user should see the system's effect
    (\enquote{you saved X kr this month by shifting dishwasher runs to
    off-peak hours}) without needing to understand how.
  \item[Graceful failure] If a device cannot communicate, it should
    continue operating at reasonable defaults rather than failing silently
    or requiring manual intervention.
  \item[Clear error messages] If something goes wrong, messages should be
    actionable: \enquote{the dishwasher lost contact with the controller,
    check that the controller is powered on} rather than
    \enquote{AEAD authentication failure on frame 4712}.
\end{description}

\subsection{Consequences of the design}

\begin{description}
  \item[Compliance] The system must be designed with GDPR and EU
    metering regulations in mind from the start, not as an afterthought.
  \item[Usability--security tension] Making the system easy to use may
    require relaxing some security measures (\eg allowing pairing with
    just a button press, which provides weaker authentication than a
    QR code with a long secret).
  \item[Third-party data sharing] Providing APIs for third-party services
    introduces new attack surfaces and privacy risks.
    Consent management and access-token scoping are essential.
  \item[Accessibility] The system must be usable by people with
    disabilities, which may affect the choice of pairing methods
    (not everyone can scan a QR code).
\end{description}


\section{Future work and open problems}\label{EnergyFutureWork}

Several aspects of smart energy management systems remain active areas of
research and development:

\begin{description}
  \item[Post-quantum cryptography] The symmetric algorithms (AES,
    HMAC) used in device-to-controller communication are largely
    quantum-resistant.
    However, any certificate-based or Diffie--Hellman-based key exchange
    must be migrated to post-quantum alternatives (ML-KEM, ML-DSA) to
    ensure long-term security.
  \item[Denial-of-service on the radio channel] Jamming and flooding
    attacks on the wireless channel can disrupt coordination.
    Advanced techniques such as spread-spectrum communication,
    cryptographic puzzles for rate limiting, and redundant communication
    paths (\eg combining radio with power-line communication) could
    improve resilience.
  \item[Multi-household coordination] In apartment buildings, multiple
    households share the same electrical supply.
    Coordinating across households raises privacy challenges: how can
    the building optimize collective consumption without any household
    revealing its individual usage?
    Techniques from secure multi-party computation could enable this.
  \item[Renewable energy and local trading] Households with solar panels
    may want to sell excess electricity to neighbours.
    This peer-to-peer energy trading requires authenticated and
    privacy-preserving transaction protocols.
  \item[Supply-chain security for devices] How do we ensure that smart
    devices have not been tampered with during manufacturing, transport,
    or retail?
    Hardware attestation and supply-chain verification protocols are
    relevant.
  \item[Machine-learning on encrypted data] The controller could use
    machine learning to predict consumption patterns and optimize
    scheduling.
    If this processing moves to the cloud, homomorphic encryption or
    federated learning could preserve privacy.
\end{description}

\begin{exercise}
  Choose one of the topics above and outline how you would approach
  designing a solution.
  What cryptographic primitives would you use?
  What are the main challenges?
\end{exercise}


\section*{Acknowledgements}

This chapter was written by Daniel Bosk with the assistance of Dan-Claude
van Thropic.
